\documentclass[11pt]{article}
\usepackage[margin=1in]{geometry}
\usepackage{array}
\usepackage{enumitem}
\usepackage{xcolor}
\usepackage{hyperref}
\usepackage{listings}

\hypersetup{colorlinks=true, linkcolor=blue, urlcolor=blue}
\setlength{\parindent}{0pt}
\setlength{\parskip}{6pt}
\setlist[itemize]{topsep=3pt,itemsep=2pt,leftmargin=*}
\setlist[enumerate]{topsep=3pt,itemsep=4pt,leftmargin=*}

\lstset{
  basicstyle=\ttfamily\small,
  breaklines=true,
  frame=single,
  columns=fullflexible,
  keepspaces=true,
  showstringspaces=false,
  backgroundcolor=\color{gray!6}
}

\definecolor{predictcolor}{RGB}{255,242,200}
\definecolor{discusscolor}{RGB}{214,234,255}
\definecolor{funfactcolor}{RGB}{222,247,222}
\definecolor{debatecolor}{RGB}{255,221,221}

\newcommand{\calloutbox}[3]{%
  \vspace{4pt}\noindent
  \colorbox{#1}{\parbox{0.96\linewidth}{\textbf{#2}}}\\[2pt]
  \noindent\fbox{\begin{minipage}{0.94\linewidth}\vspace{2pt}#3\vspace{2pt}\end{minipage}}
  \vspace{6pt}
}
\newcommand{\predictbox}[1]{\calloutbox{predictcolor}{PREDICTION TIME}{#1}}
\newcommand{\discussbox}[1]{\calloutbox{discusscolor}{HUDDLE UP}{#1}}
\newcommand{\funfact}[1]{\calloutbox{funfactcolor}{FUN FACT}{#1}}
\newcommand{\debatebox}[1]{\calloutbox{debatecolor}{DEBATE CORNER}{#1}}
\newcommand{\claimbox}[1]{\calloutbox{funfactcolor}{CLAIM, EVIDENCE, CONTEXT}{#1}}
\newcommand{\badge}[1]{$[\ ]$ \textbf{#1}}

\newcommand{\writebox}[1]{
  \vspace{3pt}
  \noindent\fbox{
    \begin{minipage}{0.96\linewidth}
      \textbf{#1}\\[12pt]
      \rule{\linewidth}{0.4pt}\\[14pt]
      \rule{\linewidth}{0.4pt}\\[14pt]
      \rule{\linewidth}{0.4pt}\\[14pt]
    \end{minipage}
  }
  \vspace{6pt}
}

\begin{document}

\begin{center}
  {\LARGE MA 213 Week 1: Lab 1}\\[3pt]
  {\large R Intro and first data exploration}
\end{center}

\begin{enumerate}
  \item \textbf{Your name:} \underline{\hspace{0.3\textwidth}} \hfill \textbf{BU ID:} \underline{\hspace{0.25\textwidth}}
  \item \textbf{Partner's name:} \underline{\hspace{0.3\textwidth}} \hfill \textbf{BU ID:} \underline{\hspace{0.25\textwidth}}
\end{enumerate}

\textbf{Big idea:} In this lab, you will practice using R to explore data sets and answer questions about them. You will also practice writing up your findings in a clear and concise way.
\textbf{Do not use GenAI tools for this lab.} The point is to practice your own analysis work and coding skills. 

\textbf{File used:} \texttt{sleep.csv}

\textbf{Your mission:} Work with your partner to get familiar with R, import a data set, and figure out what is hidden in the structure of the data.

\section*{Icebreaker}

Before opening R: introduce yourself to your partner and share one fun fact about yourself.

\writebox{One thing I want my partner to know about me:}

\section*{Question 1. Import the data and explore it}

Import \texttt{sleep.csv} into R. Then look at the first few rows.

\begin{lstlisting}[language=R]
# Read the CSV file.
sleep_data <- read.csv("sleep.csv", row.names = 1)

# Look at the first few rows.
head(sleep_data)
\end{lstlisting}

\discussbox{Before doing calculations, discuss with your partner: What do you think this data set is about? What do you expect to find?}

\writebox{Write 2--3 sentences describing what you see in the data.}

\section*{Question 2. How many patients are in the data set? How many variables? What are the variable names?}

\begin{lstlisting}[language=R]
# Hint: use dim() and names() to explore the data.
# The sleep data is also a built-in R data set, so help(sleep) may help.
# https://stat.ethz.ch/R-manual/R-devel/library/datasets/html/sleep.html

dim(sleep_data)
names(sleep_data)
length(unique(sleep_data$ID))
\end{lstlisting}

\discussbox{If R gives an error such as ``cannot open file,'' it usually means R is looking in a different folder. Use \texttt{getwd()} to see your working directory.}

\writebox{Number of rows: \hspace{2cm} Number of variables: \hspace{2cm} Number of patients: \hspace{2cm}}

\section*{Question 3. What does one row represent?}

Use \texttt{str()} to inspect the variables. Then decide which variables are numerical and which are categorical.

\begin{lstlisting}[language=R]
str(sleep_data)
\end{lstlisting}

\begin{center}
\begin{tabular}{|p{0.22\textwidth}|p{0.28\textwidth}|p{0.4\textwidth}|}
\hline
\textbf{Variable} & \textbf{Categorical or numerical?} & \textbf{What does it mean?}\\
\hline
\texttt{extra} & & \\[22pt]
\hline
\texttt{group} & & \\[22pt]
\hline
\texttt{ID} & & \\[22pt]
\hline
\end{tabular}
\end{center}

\writebox{What does one row represent in this data set?}

\section*{Question 4. What is hidden in the data set?}

At first, the data set may look like it has 20 different people. Look more carefully at \texttt{ID}. Is that true?

\begin{lstlisting}[language=R]
table(sleep_data$ID)
table(sleep_data$group)

# This shows each patient's records together.
sleep_data[order(sleep_data$ID, sleep_data$group), ]
\end{lstlisting}

\discussbox{The hidden structure is that each patient appears twice: once in group 1 and once in group 2. This means the data are paired.}

\begin{center}
\begin{tabular}{|p{0.3\textwidth}|p{0.58\textwidth}|}
\hline
\textbf{Question} & \textbf{Your answer}\\
\hline
How many rows are there? & \\[18pt]
\hline
How many unique patients are there? & \\[18pt]
\hline
How many times does each patient appear? & \\[18pt]
\hline
What is hidden in the data? & \\[24pt]
\hline
\end{tabular}
\end{center}

\section*{Question 5. Compare the two groups}

Now compare the amount of extra sleep in group 1 and group 2.

\begin{lstlisting}[language=R]
aggregate(extra ~ group, data = sleep_data, FUN = mean)
aggregate(extra ~ group, data = sleep_data, FUN = median)
\end{lstlisting}

\writebox{Which group had the larger average extra sleep? How do you know?}

\claimbox{Final answer: The hidden feature of the data set is \underline{\hspace{5cm}}. The group with more extra sleep is \underline{\hspace{3cm}}. My evidence is \underline{\hspace{5cm}}.}

\textbf{Submission check:} Make sure your R script runs from top to bottom and includes your code for Questions 1--5.

\end{document}
